\documentclass[12pt]{article}
\usepackage[utf8]{inputenc}
\usepackage[T1]{fontenc}
\usepackage[margin=1in]{geometry}
\usepackage{setspace}
\usepackage{xcolor}
\usepackage[colorlinks=true, urlcolor=blue!60!black, linkcolor=blue!60!black]{hyperref}
\setlength{\parindent}{2em}
\setlength{\parskip}{1.25em}
\renewcommand{\baselinestretch}{1.0}
\title{\textit{Pike v. Bruce Church, Inc.}, 397 U.S. 137 (1970)}
\author{Case Brief}
\date{}
\begin{document}
\maketitle

\section{Case Metadata}

\begin{description}
  \item[Case Name:] \textit{Pike v. Bruce Church, Inc.} (Loren J. Pike, Supervisor of Inspection, Arizona Fruit and Vegetable Standardization Act, Appellant, v. Bruce Church, Inc., Appellee)
  \item[Citation:] 397 U.S. 137 (1970); 90 S. Ct. 844
  \item[Docket No.:] No. 301
  \item[Court:] Supreme Court of the United States (direct appeal under 28 U.S.C. \S~1253 from a three-judge U.S. District Court for the District of Arizona convened under 28 U.S.C. \S\S~2281, 2284)
  \item[Argued:] January 13, 1970
  \item[Decided:] March 2, 1970
  \item[Vote:] 8--0 (unanimous; one vacancy existed between Justice Abe Fortas's resignation and Justice Harry Blackmun's appointment)
  \item[Majority Opinion Author:] Justice Potter Stewart (joined by Chief Justice Warren E. Burger and Justices Hugo L. Black, William O. Douglas, John Marshall Harlan II, William J. Brennan, Jr., Byron R. White, and Thurgood Marshall)
  \item[Concurring / Dissenting Opinions:] None
  \item[Constitutional Provision(s):] U.S. Const. art. I, \S~8, cl. 3 (Commerce Clause; Dormant Commerce Clause)
  \item[Chapter:] Federalism and State Power
  \item[Primary Opinion Sources:] \href{https://supreme.justia.com/cases/federal/us/397/137/}{Justia U.S. Supreme Court Center} $|$ \href{https://www.law.cornell.edu/supremecourt/text/397/137}{Cornell Legal Information Institute (LII)}
\end{description}

\section{Facts and Procedural History}

\subsection{Factual Background}

Bruce Church, Inc. (appellee) conducted extensive commercial farming operations---growing, harvesting, processing, and packing fruits and vegetables---in Arizona and California for interstate shipment nationwide (\href{https://supreme.justia.com/cases/federal/us/397/137/}{397 U.S. at 138--139}). Pursuant to a 1964 lease with the Secretary of the Interior, the Colorado River Indian Agency, and the Colorado River Indian Tribes, the company undertook to develop approximately 6,400 acres of arid, uncultivated land in Parker, Arizona, investing more than \$3,000,000 to clear, level, irrigate, and develop the land (\href{https://supreme.justia.com/cases/federal/us/397/137/}{397 U.S. at 139}).

Starting in 1966, Bruce Church harvested large annual crops of Parker cantaloupes regarded as higher in quality than those grown elsewhere in Arizona (\href{https://supreme.justia.com/cases/federal/us/397/137/}{397 U.S. at 139}). Because cantaloupes are highly perishable, they must be sorted, inspected, packed, and shipped immediately upon maturity to prevent spoilage, and commercial-scale packing can be performed only in specialized packing sheds (\href{https://supreme.justia.com/cases/federal/us/397/137/}{397 U.S. at 139--140}). Lacking packing facilities at Parker, Bruce Church transported its 1966 and 1967 harvests in bulk loads 31 miles across the state line to its centralized, efficient packing sheds in Blythe, California (\href{https://supreme.justia.com/cases/federal/us/397/137/}{397 U.S. at 139}). There, the melons were packed in standard containers under California's Fruit, Nut and Vegetable Standardization Act (Cal. Agric. Code, Div. 17, \S~45691)---stipulated to be virtually identical to Arizona's statute---and shipped in crates bearing the California packer's name rather than identifying the fruit as Arizona-grown (\href{https://supreme.justia.com/cases/federal/us/397/137/}{397 U.S. at 144 \& n.7}).

In 1968, Loren J. Pike (appellant), the official enforcing the Arizona Fruit and Vegetable Standardization Act (Ariz. Rev. Stat. Ann., Tit. 3, c. 3, Art. 4), invoked Ariz. Rev. Stat. Ann. \S~3-503C (Supp. 1969), requiring all Arizona-grown cantaloupes offered for sale to ``be packed in regular compact arrangement in closed standard containers approved by the supervisor'' (\href{https://supreme.justia.com/cases/federal/us/397/137/}{397 U.S. at 138 \& nn.1--2}). Pike ordered Bruce Church not to ship uncrated cantaloupes from Parker to Blythe unless first packed in approved containers inside Arizona (\href{https://supreme.justia.com/cases/federal/us/397/137/}{397 U.S. at 138--139}). Lacking any nearby Arizona packing shed, Bruce Church faced the imminent, unrecoverable loss of its anticipated 1968 crop valued at \$700,000 gross (\href{https://supreme.justia.com/cases/federal/us/397/137/}{397 U.S. at 139--140}).

\subsection{Procedural History}

Bruce Church sued in federal court to enjoin Pike's order as an unconstitutional burden on interstate commerce (\href{https://supreme.justia.com/cases/federal/us/397/137/}{397 U.S. at 138}). A three-judge District Court (28 U.S.C. \S\S~2281, 2284) granted preliminary relief to prevent the \$700,000 crop loss, properly declining to abstain pending state-court litigation given the emergency and narrow as-applied challenge (\href{https://supreme.justia.com/cases/federal/us/397/137/}{397 U.S. at 138, 140 \& n.3}). After discovery, the parties stipulated that Pike's order would compel Bruce Church to build and operate packing facilities in or near Parker taking many months and costing approximately \$200,000 (\href{https://supreme.justia.com/cases/federal/us/397/137/}{397 U.S. at 140}). The three-judge District Court issued a permanent injunction (unreported), holding that the order unlawfully burdened interstate commerce (\href{https://supreme.justia.com/cases/federal/us/397/137/}{397 U.S. at 138, 140 \& n.4}). Pike took a direct appeal to the Supreme Court under 28 U.S.C. \S~1253 (396 U.S. 812) (\href{https://supreme.justia.com/cases/federal/us/397/137/}{397 U.S. at 138}).

\section{Issue(s) Presented}

\begin{enumerate}
  \item \textbf{Threshold Interstate Commerce Nexus:} Does a state administrative order barring the out-of-state bulk shipment of locally grown produce unless first packed in-state burden ``interstate commerce,'' or does it regulate merely intrastate preparation before interstate commerce begins?
  \item \textbf{Dormant Commerce Clause Validity:} May a state enforce a facially neutral agricultural standardization statute to prohibit a grower from shipping superior-quality produce 31 miles across a state border for packing---forcing an unneeded \$200,000 in-state capital expenditure---solely to enhance the reputation of other in-state growers by identifying the produce as originating in the state?
\end{enumerate}

\section{Holding}

\textbf{Yes} as to the threshold burden, and \textbf{No} as to constitutionality; judgment \textbf{affirmed} (8--0) (\href{https://supreme.justia.com/cases/federal/us/397/137/}{397 U.S. at 140--146}).

\begin{enumerate}
  \item Pike's order burdens interstate commerce because the perishable cantaloupes were destined for immediate shipment to an ascertainable out-of-state location upon harvest, and the order required an operation conducted out-of-state to be performed in-state so Arizona could regulate it (\href{https://supreme.justia.com/cases/federal/us/397/137/}{397 U.S. at 140--142}).
  \item Where a statute regulates evenhandedly to effectuate a legitimate local public interest and its effects on interstate commerce are only incidental, it stands unless the burden on interstate commerce is ``clearly excessive in relation to the putative local benefits'' (\href{https://supreme.justia.com/cases/federal/us/397/137/}{397 U.S. at 142}). Arizona's minimal interest in identifying the origin of Bruce Church's superior cantaloupes to enhance other growers' reputations cannot justify forcing the company to build and operate an unneeded \$200,000 packing plant in Arizona (\href{https://supreme.justia.com/cases/federal/us/397/137/}{397 U.S. at 145--146}).
\end{enumerate}

\section{Reasoning}

\subsection{Threshold Rejection of the ``Intrastate Preparation'' Defense}

Justice Stewart rejected Pike's claim that the order regulated only intrastate packing prior to interstate shipment (\href{https://supreme.justia.com/cases/federal/us/397/137/}{397 U.S. at 140--141}). Distinguishing \textit{Federal Compress \& Warehouse Co. v. McLean}, 291 U.S. 17 (1934), and \textit{Chassaniol v. City of Greenwood}, 291 U.S. 584 (1934), the Court noted that the taxed cotton there had come to rest in Mississippi with ``no ascertainable destination without the state,'' whereas Bruce Church's cantaloupes were destined immediately for Blythe, California (\href{https://supreme.justia.com/cases/federal/us/397/137/}{397 U.S. at 141}). More fundamentally, Mississippi taxed existing local operations, whereas Arizona demanded that an out-of-state operation be relocated in-state for local regulation---a forced-processing mandate consistently invalidated in \textit{Foster-Fountain Packing Co. v. Haydel}, 278 U.S. 1 (1928), \textit{Johnson v. Haydel}, 278 U.S. 16 (1928), and \href{https://supreme.justia.com/cases/federal/us/334/385/}{\textit{Toomer v. Witsell}, 334 U.S. 385 (1948)} (\href{https://supreme.justia.com/cases/federal/us/397/137/}{397 U.S. at 141--142}).

\subsection{The \textit{Pike} Balancing Framework}

Synthesizing \textit{Huron Portland Cement Co. v. Detroit}, 362 U.S. 440, 443 (1960), and \href{https://supreme.justia.com/cases/federal/us/325/761/}{\textit{Southern Pacific Co. v. Arizona}, 325 U.S. 761 (1945)}, the Court articulated a three-part inquiry for facially neutral laws (\href{https://supreme.justia.com/cases/federal/us/397/137/}{397 U.S. at 142}): (1) whether the law regulates evenhandedly for a legitimate local public interest with only incidental interstate effects; (2) whether the interstate burden is ``clearly excessive in relation to the putative local benefits''; and (3) how the nature of the local interest weighs against whether it ``could be promoted as well with a lesser impact on interstate activities.''

\subsection{Minimal Local Interest vs. Forced In-State Operations}

While Arizona's 1929 Act generally served legitimate ends---preventing deceptive packaging, barring contaminated produce (\textit{Sligh v. Kirkwood}, 237 U.S. 52 (1915)), and regulating local containers (\textit{Pacific States Box \& Basket Co. v. White}, 296 U.S. 176 (1935))---Pike's as-applied order did not involve health, safety, or consumer deception (\href{https://supreme.justia.com/cases/federal/us/397/137/}{397 U.S. at 142--144}). Bruce Church's superior cantaloupes were packed under California's matching standards and were not labeled as Arizona produce; Arizona sought solely to capture ``reflected good will'' for other Arizona growers---an interest the Court deemed ``tenuous'' and ``minimal, at best'' (\href{https://supreme.justia.com/cases/federal/us/397/137/}{397 U.S. at 144--146}).

On the burden side, ``[t]he nature of that burden is, constitutionally, more significant than its extent'' (\href{https://supreme.justia.com/cases/federal/us/397/137/}{397 U.S. at 145}). Laws requiring business operations to be performed in-state when they can be performed more efficiently elsewhere have been declared ``virtually per se illegal,'' even absent a concealed purpose to hoard local employment (\href{https://supreme.justia.com/cases/federal/us/397/137/}{397 U.S. at 145--146}, quoting \href{https://supreme.justia.com/cases/federal/us/334/385/}{\textit{Toomer}, 334 U.S. at 403--404}). Because the Commerce Clause forbids mandating local processing to secure local employment, \textit{a fortiori} it forbids placing a ``straitjacket'' on Bruce Church's interstate resources solely to boost other growers' reputations (\href{https://supreme.justia.com/cases/federal/us/397/137/}{397 U.S. at 146}).

\section{Concurring and Dissenting Opinions}

None. Justice Stewart delivered the unanimous opinion of an eight-member Court (\href{https://supreme.justia.com/cases/federal/us/397/137/}{397 U.S. at 138}).

\section{Key Quotations}

\begin{enumerate}
  \item \textbf{The Canonical \textit{Pike} Balancing Test} (\href{https://supreme.justia.com/cases/federal/us/397/137/}{397 U.S. at 142}; \href{https://www.law.cornell.edu/supremecourt/text/397/137}{Cornell LII}):
  \begin{quote}
    ``Although the criteria for determining the validity of state statutes affecting interstate commerce have been variously stated, the general rule that emerges can be phrased as follows: where the statute regulates evenhandedly to effectuate a legitimate local public interest, and its effects on interstate commerce are only incidental, it will be upheld unless the burden imposed on such commerce is clearly excessive in relation to the putative local benefits. \textit{Huron Cement Co. v. Detroit}, 362 U. S. 440, 362 U. S. 443. If a legitimate local purpose is found, then the question becomes one of degree. And the extent of the burden that will be tolerated will, of course, depend on the nature of the local interest involved, and on whether it could be promoted as well with a lesser impact on interstate activities.''
  \end{quote}

  \item \textbf{Unmasking Arizona's Reputational Objective} (\href{https://supreme.justia.com/cases/federal/us/397/137/}{397 U.S. at 144}; \href{https://www.law.cornell.edu/supremecourt/text/397/137}{Cornell LII}):
  \begin{quote}
    ``The appellant's order would forbid the company to pack its cantaloupes outside Arizona not for the purpose of keeping the reputation of its growers unsullied, but to enhance their reputation through the reflected good will of the company's superior produce. The appellant, in other words, is not complaining because the company is putting the good name of Arizona on an inferior or deceptively packaged product, but because it is not putting that name on a product that is superior and well packaged.''
  \end{quote}

  \item \textbf{Suspicion of Forced Local-Processing Mandates} (\href{https://supreme.justia.com/cases/federal/us/397/137/}{397 U.S. at 145}; \href{https://www.law.cornell.edu/supremecourt/text/397/137}{Cornell LII}):
  \begin{quote}
    ``The nature of that burden is, constitutionally, more significant than its extent. For the Court has viewed with particular suspicion state statutes requiring business operations to be performed in the home State that could more efficiently be performed elsewhere. Even where the State is pursuing a clearly legitimate local interest, this particular burden on commerce has been declared to be virtually per se illegal.''
  \end{quote}

  \item \textbf{The ``Straitjacket'' and Final \textit{A Fortiori} Conclusion} (\href{https://supreme.justia.com/cases/federal/us/397/137/}{397 U.S. at 146}; \href{https://www.law.cornell.edu/supremecourt/text/397/137}{Cornell LII}):
  \begin{quote}
    ``While the order issued under the Arizona statute does not impose such rigidity on an entire industry, it does impose just such a straitjacket on the appellee company with respect to the allocation of its interstate resources. Such an incidental consequence of a regulatory scheme could perhaps be tolerated if a more compelling state interest were involved. But here, the State's interest is minimal, at best -- certainly less substantial than a State's interest in securing employment for its people. If the Commerce Clause forbids a State to require work to be done within its jurisdiction to promote local employment, then surely it cannot permit a State to require a person to go into a local packing business solely for the sake of enhancing the reputation of other producers within its borders.''
  \end{quote}
\end{enumerate}

\section{Doctrinal Significance and Modern Legacy (Through 2026)}

\subsection{Two-Tier Framework and Pre-2023 Evolution}
Paired with \href{https://supreme.justia.com/cases/federal/us/437/617/}{\textit{City of Philadelphia v. New Jersey}, 437 U.S. 617 (1978)}, \textit{Pike} cemented the two-tier Dormant Commerce Clause: discriminatory laws face ``virtually per se'' invalidity, while evenhanded laws with incidental effects undergo \textit{Pike} balancing. In practice, the Court rarely struck down nondiscriminatory regulations, upholding market-structure and packaging rules in \href{https://supreme.justia.com/cases/federal/us/437/117/}{\textit{Exxon Corp. v. Governor of Maryland}, 437 U.S. 117, 127 (1978)} (protecting the interstate market, not particular firms' ``methods of operation''), and \href{https://supreme.justia.com/cases/federal/us/449/456/}{\textit{Minnesota v. Clover Leaf Creamery Co.}, 449 U.S. 456 (1981)}. Rare invalidations---such as \href{https://supreme.justia.com/cases/federal/us/450/662/}{\textit{Kassel v. Consolidated Freightways Corp.}, 450 U.S. 662 (1981)}, and \href{https://supreme.justia.com/cases/federal/us/457/624/}{\textit{Edgar v. MITE Corp.}, 457 U.S. 624 (1982)}---involved interstate transportation or direct out-of-state tender-offer regulation. Meanwhile, Justice Scalia attacked \textit{Pike} as weighing ``whether a particular line is longer than a particular rock is heavy'' (\href{https://supreme.justia.com/cases/federal/us/486/888/}{\textit{Bendix Autolite Corp. v. Midwesco Enters.}, 486 U.S. 888, 897 (1988)}), and \href{https://www.law.cornell.edu/supremecourt/text/06-666}{\textit{Dept. of Revenue of Ky. v. Davis}, 553 U.S. 328, 353--355 (2008)}, warned of judicial institutional limits.

\subsection{\textit{National Pork Producers Council v. Ross} (2023) and Developments Through 2026}
In \href{https://www.law.cornell.edu/supremecourt/text/21-468}{\textit{National Pork Producers Council v. Ross}, 598 U.S. 356 (2023)}, upholding California's Proposition 12 pork-confinement sales ban, the Court unanimously rejected any ``almost per se'' extraterritoriality rule (Part III) and fractured over \textit{Pike}:
\begin{itemize}
  \item \textbf{Heartland Reframing (5--4, Part IV-A):} Justice Gorsuch (joined by Thomas, Sotomayor, Kagan, and Barrett) recast \textit{Pike} as primarily a tool to ``smoke out'' latent protectionism---reinterpreting \textit{Pike} itself as insulating local packers---or protect transportation instrumentalities.
  \item \textbf{Incommensurability (3 Justices, Parts IV-B \& IV-D):} Gorsuch, Thomas, and Barrett argued courts cannot weigh economic costs against moral or health interests.
  \item \textbf{Threshold Burden Failure (4 Justices, Part IV-C):} Gorsuch, Thomas, Sotomayor, and Kagan held under \href{https://supreme.justia.com/cases/federal/us/437/117/}{\textit{Exxon}} that petitioners failed to plead a substantial market-wide burden.
  \item \textbf{Six-Justice Retention:} Chief Justice Roberts (joined by Alito, Kavanaugh, and Jackson) dissented on \textit{Pike}, noting that six Justices (adding Sotomayor and Kagan) preserved \textit{Pike} balancing and five (adding Barrett) saw a plausibly alleged burden.
\end{itemize}

Through 2026, lower courts apply \textit{Ross} strictly: in \href{https://law.justia.com/cases/federal/appellate-courts/ca1/24-1759/24-1759-2025-10-03.html}{\textit{Triumph Foods, LLC v. Campbell} (1st Cir. 2025)}, \href{https://www.scotusblog.com/cases/triumph-foods-llc-v-campbell/}{cert. denied, No. 25-1047 (June 30, 2026)}, the First Circuit upheld Massachusetts's Question 3 pork law, and in \href{https://www.bdlaw.com/publications/federal-court-upholds-oregon-packaging-epr-law-providing-important-early-roadmap-for-epr-challenges/}{\textit{Nat'l Ass'n of Wholesaler-Distributors v. Feldon}, No. 3:25-cv-01334-SI (D. Or. Aug. 27, 2026)}, a district court upheld Oregon's packaging EPR statute after a bench trial, holding that steep distributor fees did not show discrimination or a market-wide burden.

\section{Pedagogical Critiques and Exam Hypotheticals}

\subsection{Pedagogical Critiques}
\begin{enumerate}
  \item \textbf{Hidden Protectionism vs. True Balancing:} Although famed for balancing, \textit{Pike} relied heavily on the ``virtually per se illegal'' rule against forced local processing (\href{https://supreme.justia.com/cases/federal/us/397/137/}{397 U.S. at 145}).
  \item \textbf{Incommensurability and Unexplored Alternatives:} \textit{Pike} offers no metric for weighing compliance dollars against noneconomic values (\href{https://supreme.justia.com/cases/federal/us/486/888/}{\textit{Bendix}}, 486 U.S. at 897) and omitted an obvious less restrictive alternative---requiring origin stamps on California-packed crates.
\end{enumerate}

\subsection{Exam Hypotheticals}
\textbf{1. Forced Local Operations (\textit{Pike} Violated):} State A requires all locally grown spinach to be logged on physical servers inside State A before interstate shipment to aid health inspections, forcing a border grower to build a redundant \$450,000 in-state data center despite instantaneous remote access from its State B facility. Even post-\href{https://www.law.cornell.edu/supremecourt/text/21-468}{\textit{Ross}}, mandating local operations when remote access achieves the same goal fails \textit{Pike}.

\textbf{2. ``Opposite Side'' Contrast (State Prevails):} State C bans retail sale of produce in single-use polystyrene clamshells to curb landfill waste, applying identically to in-state and out-of-state growers with no local-packing mandate. An out-of-state grower faces \$1.2 million in retooling costs. Under \href{https://supreme.justia.com/cases/federal/us/449/456/}{\textit{Clover Leaf Creamery}}, \href{https://supreme.justia.com/cases/federal/us/437/117/}{\textit{Exxon}}, and \href{https://www.law.cornell.edu/supremecourt/text/21-468}{\textit{Ross}}, State C wins: there is no forced in-state processing, firm-level retooling is merely a shift in ``methods of operation,'' and the environmental benefit is substantial.

\end{document}
